\section{平方完成}\label{節:平方完成と積分}

前節では,分母が一次式の積に分かれていた.では,分母や根号の中に二次式 $x^2+bx+c$ があり,因数分解がすぐには見えない場合はどうすればよいだろうか.ここで1つの知識を思い出そう.平方完成である.
\[
 x^2+bx+c=\left(x+\frac b2\right)^2+c-\frac{b^2}4
\]
$m=-\frac b2$,\,$k=c-\frac{b^2}4$ と置けば,この二次式は
\[
 (x-m)^2+k
\]
と書ける.\,$x$ そのものよりも,\,$x-m$ を1つのまとまりとして見ると,二次式の中にある構造が明確になる.本節では,このまとまりを手がかりに,前節の部分分数か,合成関数の微分の逆に帰着させる.

進み方は $k$ の符号で分かれる.\,$k<0$ のときは,\,$k=-h^2$ $(h>0)$ と書けば平方の差になり,
\[
 (x-m)^2-h^2=(x-m-h)(x-m+h)
\]
と因数分解できる.あとは前節の部分分数に帰着する.\,$k=0$ のときは $(x-m)^2$ であり,一次式のべきとして積分できる.

$k>0$ のときは,\,$(x-m)^2+k\ge k>0$ なので,実数の一次式の積には分解できない.このときは,まとまり $(x-m)^2+k$ を微分すると $2(x-m)$ になることを使う.積分される式の中に $x-m$ が掛かっていれば,合成関数の微分を逆に読んで原始関数の候補を作れる.

\begin{bookpoint}
\textbf{POINT}\par
$x^2+bx+c=(x-m)^2+k$ と平方完成し,\,$k<0$ なら因数分解して部分分数へ,\,$k>0$ なら $x-m$ が掛かっているかを見て合成関数の微分を逆に読む.
\end{bookpoint}

\begin{example}{根号の微分と積分}
$\int\frac{x-2}{\sqrt{x^2-4x+5}}\dd x$ を求めよ.
\end{example}

\paragraphbf{方針}
根号の中を $(x-2)^2+1$ と平方完成し,\,$\sqrt{(x-2)^2+1}$ を微分した式と見比べる.

\paragraphbf{解答}
$\sqrt{x^2-4x+5}+C$

\paragraphbf{解法}
平方完成すると,積分される式は
\[
 \frac{x-2}{\sqrt{(x-2)^2+1}}
\]
となる.分子に $x-2$ があり,分母に平方根があるので,\,$\sqrt{(x-2)^2+1}$ を微分してみよう.合成関数の微分から
\begin{align*}
 \frac{\dd}{\dd x}\sqrt{(x-2)^2+1}
 &=\frac1{2\sqrt{(x-2)^2+1}}\cdot2(x-2)\\
 &=\frac{x-2}{\sqrt{(x-2)^2+1}}
\end{align*}
となる.これがちょうど元の式なので,
\[
 \int\frac{x-2}{\sqrt{x^2-4x+5}}\dd x
 =\sqrt{x^2-4x+5}+C
\]
を得る.根号の中は $(x-2)^2+1\ge1$ なので,この答えは実数全体で使える.

\paragraphbf{まとまりの微分がない場合}
同じ根号でも,分子に $x-2$ がない $\int\frac1{\sqrt{(x-2)^2+1}}\dd x$ では,いまの微分の結果と一致しない.その違いを無視して同じ答えにはできない.また,\,$k>0$ で分母が $(x-m)^2+k$ そのものである $\int\frac1{(x-m)^2+k}\dd x$ も,この節の変形だけでは終わらない.前者は\sectionref{節:平方根と部分積分}で,後者は\sectionref{節:三角置換の初歩}の三角置換で扱う.平方完成はいつでも積分を終わらせる操作ではないが,その後に因数分解を使えるのか,置換を考えるべきなのかを判断する材料になるのだ.
